\documentclass[10pt]{article}
\usepackage[a4paper,margin=1in]{geometry}
\usepackage[T1]{fontenc}
\usepackage{lmodern}
\usepackage{amsmath,amssymb,mathtools,bm}
\usepackage{algorithm}
\usepackage[noend]{algpseudocode}
\usepackage{microtype}
\usepackage{hyperref}

\algrenewcommand\algorithmicrequire{\textbf{Input:}}
\algrenewcommand\algorithmicensure{\textbf{Output:}}

\newcommand{\E}{\mathbb{E}}
\newcommand{\ind}{\mathbf{1}}
\newcommand{\clip}{\operatorname{clip}}
\newcommand{\sg}{\operatorname{sg}}
\newcommand{\norm}{\operatorname{Norm}}
\newcommand{\uid}{\operatorname{uid}}
\newcommand{\traj}{\operatorname{traj}}
\newcommand{\step}{\operatorname{step}}

\title{Method: Shared-Backbone Step-wise PPO with Episode-Level Relative Correction}
\author{}
\date{}

\begin{document}
\maketitle
\vspace{-1.5em}

\section*{Method}

We consider long-horizon agentic reinforcement learning where an episode consists of multiple environment interaction steps. At each step, the agent observes a context, generates an action such as a tool call, and receives an environment reward. Our goal is to construct a step-wise actor-critic algorithm while preserving the episode-level relative training signal used by GRPO/GiGPO-style grouped rollouts. The key design choice is to avoid maintaining a separate critic model. Instead, we attach a value head to the actor model and use the shared actor backbone to predict state values.

\subsection*{Problem formulation}

Let $x$ denote an initial task instance sampled from a training distribution. A rollout trajectory generated by policy $\pi_\theta$ is
\[
\tau = (s_0,a_0,r_0,s_1,a_1,r_1,\ldots,s_{T-1},a_{T-1},r_{T-1}),
\]
where $s_t$ is the agent context at environment step $t$, $a_t$ is the complete action generated by the language model, and $r_t\in\mathbb{R}$ is the immediate environment reward. Although $a_t$ is represented by a sequence of response tokens,
\[
a_t = (y_{t,1},y_{t,2},\ldots,y_{t,L_t}),
\]
we treat the whole response as a single environment-level action for value estimation, return computation, and advantage estimation.

For grouped rollouts, each initial task instance $x$ is repeated $K$ times. The resulting trajectories share the same episode-group identifier $g=\uid(x)$, while each trajectory has a unique trajectory identifier $m=\traj(\tau)$. We denote the set of trajectories in group $g$ by $\mathcal{T}_g$.

\subsection*{Shared actor-value model}

The actor is a causal language model with parameters $\theta$. Instead of creating a standalone critic backbone, we attach a regression head $v_\psi$ to the actor backbone. Given the hidden states of the actor model,
\[
H_t = f_\theta(s_t,a_t) \in \mathbb{R}^{n_t\times d},
\]
where $n_t$ is the full context-plus-response sequence length, we define the strict state-value estimate using the hidden state at the action-generation boundary:
\[
V_{\theta,\psi}(s_t)
= v_\psi\bigl(H_t[b_t]\bigr),
\qquad
b_t = n_t - L_t - 1.
\]
Here $b_t$ indexes the final context token before the first generated action token. Because the actor is causal, $H_t[b_t]$ does not attend to future response tokens. Therefore, $V_{\theta,\psi}(s_t)$ is a state-value estimate rather than an action-conditioned value estimate.

The value head is implemented as a small MLP,
\[
v_\psi(h) = W_2\,\sigma(W_1 h + b_1) + b_2,
\]
where $h\in\mathbb{R}^d$, $W_1\in\mathbb{R}^{d_h\times d}$, $W_2\in\mathbb{R}^{1\times d_h}$, and $\sigma$ is GELU. A linear head $v_\psi(h)=w^\top h+b$ is a valid special case. In the main method, gradients from the value loss are allowed to update both $\psi$ and the shared backbone $\theta$, with a small value-loss coefficient. Detaching the backbone for the value loss is reserved as an ablation.

\paragraph{Action-conditioned variants.}
One may instead pool hidden states over response tokens or use the last response token value. This yields an action-conditioned quantity $V(s_t,a_t)$ or $Q(s_t,a_t)$ rather than a strict state-value baseline. We do not use this in the main method, but retain it as an ablation.

\subsection*{Step-wise returns and GAE}

For each trajectory $\tau$, we compute step-wise temporal-difference residuals along the environment-step axis. Let $d_t\in\{0,1\}$ indicate whether step $t$ is terminal. With discount $\gamma\in[0,1]$ and GAE parameter $\lambda\in[0,1]$, define
\[
\delta_t
= r_t + \gamma (1-d_t) V_{\theta,\psi}^{\mathrm{old}}(s_{t+1})
  - V_{\theta,\psi}^{\mathrm{old}}(s_t),
\]
where the old values are evaluated before policy optimization on the collected batch. The step-wise GAE is computed backward over environment steps:
\[
A^{\mathrm{step}}_t
= \delta_t + \gamma\lambda(1-d_t) A^{\mathrm{step}}_{t+1},
\qquad
A^{\mathrm{step}}_{T}=0.
\]
The corresponding value target is
\[
R^{\mathrm{step}}_t
= A^{\mathrm{step}}_t + V_{\theta,\psi}^{\mathrm{old}}(s_t).
\]
Equivalently, when $\lambda=1$, $R^{\mathrm{step}}_t$ reduces to the Monte Carlo discounted return
\[
G_t = \sum_{k=t}^{T-1} \gamma^{k-t} r_k.
\]
All quantities above are indexed by environment step, not by response-token position.

\subsection*{Episode-level relative advantage}

Let the undiscounted episode return of trajectory $\tau_i\in\mathcal{T}_g$ be
\[
R_i^{\mathrm{epi}} = \sum_{t=0}^{T_i-1} r_{i,t}.
\]
Following GRPO-style group-relative normalization, the episode-level advantage is
\[
A_i^{\mathrm{epi}}
= \frac{R_i^{\mathrm{epi}} - \mu_g}{\sigma_g + \epsilon},
\qquad
\mu_g = \frac{1}{|\mathcal{T}_g|}\sum_{\tau_j\in\mathcal{T}_g} R_j^{\mathrm{epi}},
\]
\[
\sigma_g^2
= \frac{1}{|\mathcal{T}_g|-1}\sum_{\tau_j\in\mathcal{T}_g}
\bigl(R_j^{\mathrm{epi}}-\mu_g\bigr)^2.
\]
If $|\mathcal{T}_g|=1$, we set $\mu_g=0$ and $\sigma_g=1$ for numerical stability. This trajectory-level scalar is broadcast to every step in the trajectory:
\[
A^{\mathrm{epi}}_{i,t}=A_i^{\mathrm{epi}},
\qquad t=0,1,\ldots,T_i-1.
\]

\subsection*{Weighted correction of episode advantages}

The final advantage combines the episode-level relative signal and the step-wise actor-critic signal. Before mixing, both advantages are normalized over valid steps in the batch:
\[
\widehat{A}^{\mathrm{epi}}_t = \norm(A^{\mathrm{epi}}_t),
\qquad
\widehat{A}^{\mathrm{step}}_t = \norm(A^{\mathrm{step}}_t).
\]
We use a GiGPO-style weighted correction:
\[
A^{\mathrm{final}}_t
= \frac{\widehat{A}^{\mathrm{epi}}_t + w_{\mathrm{step}}\widehat{A}^{\mathrm{step}}_t}
       {1+w_{\mathrm{step}}},
\]
where $w_{\mathrm{step}}\ge0$ controls the strength of the step-wise critic correction. Setting $w_{\mathrm{step}}=0$ recovers pure episode-level GRPO-style training, while larger values place more weight on step-wise credit assignment.

\subsection*{Policy objective}

Let $\pi_{\theta_{\mathrm{old}}}$ be the behavior policy used to sample the current batch. For each action token $y_{t,\ell}$, define the token-level importance ratio
\[
\rho_{t,\ell}(\theta)
= \exp\left(
\log \pi_\theta(y_{t,\ell}\mid s_t,y_{t,<\ell})
-
\log \pi_{\theta_{\mathrm{old}}}(y_{t,\ell}\mid s_t,y_{t,<\ell})
\right).
\]
In the first implementation, the scalar step advantage $A^{\mathrm{final}}_t$ is broadcast to all valid response tokens in action $a_t$. The clipped PPO policy loss is
\[
\mathcal{L}_{\mathrm{pg}}(\theta)
= -\E_{t,\ell}\left[
\min\left(
\rho_{t,\ell}(\theta) A^{\mathrm{final}}_t,
\clip\bigl(\rho_{t,\ell}(\theta),1-\varepsilon,1+\varepsilon\bigr) A^{\mathrm{final}}_t
\right)
\right],
\]
where the expectation is taken over valid action tokens. To reduce length bias, the main configuration uses sequence-level averaging over action tokens. A stricter future variant will replace token-level ratios with action-level ratios,
\[
\rho_t^{\mathrm{act}}(\theta)
= \exp\left(\sum_{\ell=1}^{L_t}
\log \pi_\theta(y_{t,\ell}\mid s_t,y_{t,<\ell})
-
\log \pi_{\theta_{\mathrm{old}}}(y_{t,\ell}\mid s_t,y_{t,<\ell})
\right),
\]
or a length-normalized version of the same quantity.

\subsection*{Value objective}

The shared value head is trained with a clipped scalar value loss. Let
\[
V_t = V_{\theta,\psi}(s_t),
\qquad
V_t^{\mathrm{old}} = V_{\theta_{\mathrm{old}},\psi_{\mathrm{old}}}(s_t).
\]
The clipped value prediction is
\[
\widetilde{V}_t
= V_t^{\mathrm{old}}
+ \clip\bigl(V_t - V_t^{\mathrm{old}}, -\varepsilon_v, \varepsilon_v\bigr).
\]
The value loss is
\[
\mathcal{L}_{\mathrm{v}}(\theta,\psi)
= \E_t\left[
\max\left(
\bigl(V_t - R^{\mathrm{step}}_t\bigr)^2,
\bigl(\widetilde{V}_t - R^{\mathrm{step}}_t\bigr)^2
\right)
\right].
\]
This loss is scalar per environment step, even though the underlying language model can produce hidden states at every token position.

\subsection*{KL regularization and total loss}

We use a reference policy $\pi_{\mathrm{ref}}$ and add a token-level KL penalty to the actor loss. Let
\[
\mathcal{L}_{\mathrm{KL}}(\theta)
= \E_{t,\ell}\left[
D_{\mathrm{KL}}\bigl(
\pi_\theta(\cdot\mid s_t,y_{t,<\ell})
\,\|\,
\pi_{\mathrm{ref}}(\cdot\mid s_t,y_{t,<\ell})
\bigr)
\right].
\]
The final optimization objective is
\[
\mathcal{L}(\theta,\psi)
= \mathcal{L}_{\mathrm{pg}}(\theta)
+ c_v\mathcal{L}_{\mathrm{v}}(\theta,\psi)
+ \beta\mathcal{L}_{\mathrm{KL}}(\theta)
- c_H\mathcal{H}(\pi_\theta),
\]
where $c_v$ is the value-loss coefficient, $\beta$ is the KL coefficient, and $c_H$ is the entropy coefficient.

\begin{algorithm}[t]
\caption{Shared-Backbone Step-wise PPO with Episode-Level Relative Correction}
\label{alg:shared-step-ppo}
\begin{algorithmic}[1]
\Require Actor policy $\pi_\theta$ with value head $v_\psi$, reference policy $\pi_{\mathrm{ref}}$, grouped environments, group size $K$, discount $\gamma$, GAE parameter $\lambda$, step weight $w_{\mathrm{step}}$
\Ensure Updated actor and value-head parameters $(\theta,\psi)$
\For{each training iteration}
    \State Sample a batch of task instances and create $K$ grouped rollouts for each task
    \State Collect step samples $(s_t,a_t,r_t,\uid,\traj,\step)$ from the agent-environment loop
    \State Recompute old token log-probabilities under $\pi_{\theta_{\mathrm{old}}}$
    \State Compute old state values $V_t^{\mathrm{old}}=v_{\psi_{\mathrm{old}}}(H_t[b_t])$
    \State Compute episode returns $R_i^{\mathrm{epi}}$ for each trajectory
    \State Compute group-relative episode advantages $A_i^{\mathrm{epi}}$ and broadcast them to steps
    \For{each trajectory $\tau_i$}
        \State Sort its samples by $\step$
        \State Compute $A_t^{\mathrm{step}}$ and $R_t^{\mathrm{step}}$ by backward step-wise GAE
    \EndFor
    \State Normalize $A^{\mathrm{epi}}$ and $A^{\mathrm{step}}$ over valid steps
    \State Compute $A_t^{\mathrm{final}}=(A_t^{\mathrm{epi}}+w_{\mathrm{step}}A_t^{\mathrm{step}})/(1+w_{\mathrm{step}})$
    \State Broadcast $A_t^{\mathrm{final}}$ to valid response tokens for the PPO actor loss
    \For{each PPO epoch and mini-batch}
        \State Recompute current token log-probabilities under $\pi_\theta$
        \State Recompute current state values $V_t=v_\psi(H_t[b_t])$
        \State Compute clipped policy loss $\mathcal{L}_{\mathrm{pg}}$
        \State Compute clipped scalar value loss $\mathcal{L}_{\mathrm{v}}$ using $R_t^{\mathrm{step}}$
        \State Compute KL loss against $\pi_{\mathrm{ref}}$
        \State Update $(\theta,\psi)$ with $\mathcal{L}=\mathcal{L}_{\mathrm{pg}}+c_v\mathcal{L}_{\mathrm{v}}+\beta\mathcal{L}_{\mathrm{KL}}-c_H\mathcal{H}$
    \EndFor
\EndFor
\end{algorithmic}
\end{algorithm}

\subsection*{Implementation notes}

The value head is attached to the actor module before FSDP wrapping, so it is optimized and checkpointed together with the actor. Since inference engines such as vLLM or SGLang only consume language-model weights, parameters whose names contain \texttt{value\_head} must be filtered out before synchronizing actor weights to rollout engines. The rollout data must also store an explicit step index, because later batch reordering can otherwise destroy the environment-step order required by step-wise GAE.

\end{document}
